% =============================================================================
%  Month 1 Progress Report
%  Outlier-Aware Ensemble Selection for Robust Parkinson's Disease Prediction
%
%  Compiles with pdfLaTeX on Overleaf. No external files required.
% =============================================================================
\documentclass[11pt,a4paper]{article}

\usepackage[T1]{fontenc}
\usepackage[utf8]{inputenc}
\usepackage[margin=2.5cm]{geometry}
\usepackage{booktabs}
\usepackage{longtable}
\usepackage{array}
\usepackage{xcolor}
\usepackage{graphicx}
\usepackage{enumitem}
\usepackage{caption}
\usepackage[hidelinks]{hyperref}

\definecolor{accent}{HTML}{1F4E4C}
\definecolor{warn}{HTML}{B7770B}

\newcommand{\code}[1]{\texttt{\small #1}}
\newcommand{\status}[1]{\textbf{#1}}

\setlist[itemize]{itemsep=2pt,topsep=4pt}
\setlist[enumerate]{itemsep=2pt,topsep=4pt}

\captionsetup{font=small,labelfont=bf}

\title{\textbf{Outlier-Aware Ensemble Selection for\\ Robust Parkinson's Disease Prediction}\\[0.6em]
       \large Month 1 Progress Report}
\author{Parv\\[0.3em]
        \small Supervisor: Prof.\ K.~R.\ Sahasranand\\
        \small Department of Data Science, IIT Palakkad}
\date{July 2026}

\begin{document}
\maketitle

\begin{abstract}
\noindent
This report documents Month 1 of a four-month project on outlier-aware ensemble
selection for Parkinson's disease (PD) prediction using the Parkinson's
Progression Markers Initiative (PPMI) dataset. The objective for this month was
\emph{not} the research contribution itself, but a reproducible baseline
pipeline that Month 2 will extend. Twelve candidate models (four architectures
across three feature modalities) were trained and evaluated, together with five
baseline ensemble rules. Two distinct label-leakage problems were identified and
corrected during development; both are documented here because they materially
affect how the reported numbers should be interpreted. The pipeline is
configuration-driven, unit-tested, and fully reproducible from a single random
seed.
\end{abstract}

\vspace{0.5em}

% =============================================================================
\section{Scope and objectives}
% =============================================================================

The project proposes improving PD prediction by building an ensemble of diverse
machine learning models and applying a robust model-selection strategy inspired
by \emph{Principled Out-of-Distribution Detection via Multiple Testing}
(Magesh, Veeravalli, Roy and Jha, JMLR 2023). Rather than simple majority voting
or averaging, validation predictions are treated as structured vectors from
which ``outlier models'' can be identified and removed or down-weighted before
aggregation.

Month 1 deliberately excludes that contribution. Its objective is a clean,
reproducible baseline pipeline: data finalisation, feature engineering, a
multi-modality training pipeline, evaluation, naive ensembles, experiment
tracking, configuration, visualisation and documentation. No conformal
prediction, Benjamini--Hochberg/Yekutieli procedure, outlier detection or robust
aggregation is implemented at this stage; those belong to Month 2.

% =============================================================================
\section{Deliverable status}
% =============================================================================

\begin{longtable}{@{}p{6.6cm} p{2.6cm} p{5.4cm}@{}}
\toprule
\textbf{Deliverable} & \textbf{Status} & \textbf{Location} \\
\midrule
\endfirsthead
\toprule
\textbf{Deliverable} & \textbf{Status} & \textbf{Location} \\
\midrule
\endhead
Data finalisation pipeline            & Complete & \code{src/data/} \\
Dataset builder (patient-level, seeded, manifest) & Complete & \code{src/data/splitter.py} \\
Feature engineering --- clinical      & Complete & 23 raw columns \\
Feature engineering --- imaging       & Complete & 9 DaTscan SBR columns \\
Feature engineering --- biosignal     & \textcolor{warn}{Blocked} & code written; source not downloaded \\
Feature engineering --- combined      & Complete & \code{src/features/registry.py} \\
Training pipeline                     & Complete & \code{src/training/trainer.py} \\
Multi-modality training               & Complete for available data & 12 of a planned 16 runs \\
Evaluation suite                      & Complete & \code{src/evaluation/} \\
Baseline ensembles                    & Complete & 5 rules \\
Experiment tracking                   & Complete & \code{src/training/} \\
Configuration system                  & Complete & \code{config/config.yaml} \\
Visualisation                         & Complete$^{\dagger}$ & \code{src/visualization/plots.py} \\
Documentation                         & Complete & README, diagrams, this report \\
\bottomrule
\caption{Month 1 deliverable status. $^{\dagger}$One documented third-party
limitation; see Section~\ref{sec:limitations}.}
\label{tab:deliverables}
\end{longtable}

\noindent
Month 1 is complete for everything the available data permits. The single
outstanding functional deliverable --- the biosignal modality --- is blocked on
a dataset that has not yet been downloaded, not on engineering effort. Its code
path is written and registered; enabling it is a configuration change.

% =============================================================================
\section{Dataset and experimental design}
% =============================================================================

\subsection{Population}

The in-distribution population is drawn from the current-era PPMI Clinical
curated data cut: sporadic PD versus healthy control (HC), baseline visit only.
This yields \textbf{1{,}410 patients} (1{,}080 sporadic PD, 330 HC) from an
initial 19{,}450 rows across 212 columns.

\subsection{Splits}

Splitting is performed at the \emph{patient} level and stratified by label, with
a seeded, saved manifest recording exactly which patients fall in each split.

\begin{table}[h]
\centering
\begin{tabular}{@{}lrrl@{}}
\toprule
\textbf{Split} & \textbf{Patients} & \textbf{Positive rate} & \textbf{Purpose} \\
\midrule
Train        & 846 & 0.766 & model fitting \\
Validation   & 211 & 0.768 & model selection, ensemble weights \\
Calibration  & 211 & 0.763 & reserved for Month 2 conformal step \\
Test         & 142 & 0.768 & final reporting only \\
\bottomrule
\end{tabular}
\caption{Patient-level split composition (60/15/15/10, seed 42).}
\label{tab:splits}
\end{table}

Although the current population contains one row per patient, splitting is
implemented at patient level and asserts that no patient appears in two splits.
This keeps the guarantee correct if a future longitudinal, multi-visit
population is introduced.

\subsection{Feature modalities}

\begin{itemize}
  \item \textbf{Clinical} (23 columns): demographics plus non-motor, cognitive
        and behavioural assessments --- MoCA, HVLT, Trail Making, Symbol Digit,
        QUIP, GDS, STAI, SCOPA-AUT and MDS-UPDRS Part I.
  \item \textbf{Imaging} (9 columns): DaTscan striatal binding ratio (SBR)
        measures for caudate, putamen and striatum.
  \item \textbf{Combined} (32 columns): concatenation of the above.
  \item \textbf{Biosignal}: gait and arm-swing measures. Not yet downloaded;
        filed under \emph{Motor Assessments} in the PPMI portal rather than the
        separate \emph{Digital Sensor} category.
\end{itemize}

% =============================================================================
\section{Data leakage: two findings}
\label{sec:leakage}
% =============================================================================

Two distinct leakage problems were identified during development. Both are
reported here rather than buried, because each materially changes how the
results should be read.

\subsection{Finding 1: diagnostic circularity in clinical features}

An initial training run produced approximately 99--100\% AUC-ROC across every
model, which is implausible for this task. Investigation showed that Hoehn \&
Yahr stage and MDS-UPDRS Parts II, III and total were acting as proxies for the
diagnosis itself: in PPMI's enrolment protocol, a qualifying motor-exam finding
is part of how a participant is classified as PD, and healthy controls are
essentially always Hoehn \& Yahr stage 0. Median UPDRS-III in this data is 1 for
HC against 21 for PD.

Training on those columns therefore does not predict PD --- it recovers the
enrolment rule. They were removed. MDS-UPDRS Part I was retained: it is
patient-reported, not part of the diagnostic examination, and shows a real but
much smaller group separation.

A weaker version of the same caveat applies to imaging. PPMI's PD cohort
required an abnormal DaTscan for enrolment (a normal scan routes a participant
to the separate SWEDD cohort), so strong DaTscan performance is partly expected
rather than novel. It was retained because, unlike the motor examination,
DaTscan is an independent biological measurement rather than a restatement of
the clinical assessment.

\subsection{Finding 2: imputation computed across split boundaries}

A subsequent code review found that missing-value imputation statistics were
computed over the \emph{entire population before splitting}. Measured directly
on this dataset, the median used to fill missing BMI values was 26.213
(full population) where the correct training-only value is 26.366.

The absolute difference is small; the methodological consequence is not.
Held-out distribution information was entering the training features, so the
reported test scores were not clean held-out estimates, and the effect compounds
across every imputed column.

The data layer was restructured around an explicit contract:

\begin{enumerate}
  \item feature builders return \emph{raw} frames with missing values intact;
  \item \code{FeatureCleaner.fit()} observes \emph{training rows only},
        learning per-column medians and modes and deciding which columns receive
        a missingness indicator;
  \item \code{transform()} applies those fixed statistics unchanged to
        validation, calibration and test;
  \item the output column schema is pinned at fit time, so a category absent
        from a small split cannot silently reshape the feature matrix.
\end{enumerate}

A regression test uses a fixture whose training and held-out medians differ by
two orders of magnitude, so any reintroduction of this bug fails loudly rather
than quietly.

\subsection{Selection discipline}

Model selection and ensemble weighting use the validation split only. The test
split is evaluated once, for final reporting.

% =============================================================================
\section{Results}
% =============================================================================

\subsection{Individual models}

Train AUC is reported alongside held-out performance because several models are
visibly overfit; the training pipeline now emits a warning when the
train--test AUC gap exceeds 0.15.

\begin{table}[h]
\centering
\begin{tabular}{@{}l r r r r@{}}
\toprule
\textbf{Run} & \textbf{Train AUC} & \textbf{Test AUC} & \textbf{Test acc.} & \textbf{Test F1} \\
\midrule
combined + SVM            & 0.994 & \textbf{0.984} & 0.958 & 0.973 \\
combined + RandomForest   & 1.000 & 0.984 & 0.944 & 0.962 \\
imaging + RandomForest    & 1.000 & 0.982 & 0.937 & 0.958 \\
combined + XGBoost        & 1.000 & 0.982 & 0.944 & 0.962 \\
imaging + XGBoost         & 1.000 & 0.971 & 0.937 & 0.958 \\
imaging + NeuralNet       & 0.978 & 0.968 & 0.916 & 0.944 \\
imaging + SVM             & 0.990 & 0.965 & 0.923 & 0.949 \\
combined + NeuralNet      & 0.964 & 0.954 & 0.880 & 0.921 \\
clinical + XGBoost        & 1.000 & 0.809 & 0.824 & 0.891 \\
clinical + RandomForest   & 1.000 & 0.763 & 0.810 & 0.887 \\
clinical + SVM            & 0.919 & 0.736 & 0.761 & 0.856 \\
clinical + NeuralNet      & 0.698 & 0.707 & 0.782 & 0.875 \\
\bottomrule
\end{tabular}
\caption{Individual model performance, ranked by test ROC-AUC ($n=142$).}
\label{tab:models}
\end{table}

\subsection{Baseline ensembles}

\begin{table}[h]
\centering
\begin{tabular}{@{}l r r r r r r@{}}
\toprule
\textbf{Method} & \textbf{Acc.} & \textbf{Prec.} & \textbf{Rec.} & \textbf{F1} & \textbf{ROC-AUC} & \textbf{PR-AUC} \\
\midrule
Simple average / soft voting$^{\ast}$ & 0.951 & 0.990 & 0.945 & 0.967 & \textbf{0.9894} & 0.9972 \\
Accuracy-weighted average             & 0.951 & 0.990 & 0.945 & 0.967 & 0.9889 & 0.9970 \\
Majority vote                         & 0.944 & 0.963 & 0.963 & 0.963 & 0.9862 & 0.9945 \\
Single best model                     & 0.944 & 0.990 & 0.936 & 0.962 & 0.9817 & 0.9950 \\
\bottomrule
\end{tabular}
\caption{Baseline ensembles over all 12 models, test split ($n=142$).
$^{\ast}$Simple averaging, soft voting and equal-weighted averaging are
mathematically identical for binary classification and are reported as one row.}
\label{tab:ensembles}
\end{table}

\subsection{Interpretation}

Three observations should be carried into the eventual paper.

\begin{enumerate}
  \item \textbf{Averaging beats the single best model}, 0.9894 against 0.9817
        ROC-AUC. The ensemble premise holds, but the margin is roughly 0.008
        AUC. Month 2's method therefore faces a strong baseline, not a weak one;
        this should be stated plainly rather than presented as a large
        opportunity.
  \item \textbf{Imaging dominates clinical features}, 0.98 against 0.71--0.81
        ROC-AUC. Given the enrolment caveat in Section~\ref{sec:leakage}, this
        is partly expected rather than a novel finding.
  \item \textbf{Several models are heavily overfit}, reaching train AUC 1.000.
        This became visible only once train-split metrics were added, and should
        be addressed before Month 2 builds selection logic on these models'
        validation behaviour.
\end{enumerate}

% =============================================================================
\section{Engineering}
% =============================================================================

\subsection{Architecture}

The codebase is organised into five layers with an acyclic import graph:
foundation (\code{config}, \code{utils}), implementations (feature builders,
model wrappers), domain (registries, cleaning, splitting, evaluation),
orchestration (trainer, ensembler, visualiser) and CLI stages. No module imports
a layer above it.

Two extension points keep the architecture open for Month 2 without
modification: a new feature modality is added by subclassing
\code{FeatureBuilder} and registering it; a new architecture by subclassing
\code{ModelWrapper} and registering it. Neither requires changes to the training
loop, ensembling or visualisation, all of which iterate over whatever is
available.

\subsection{Reproducibility}

\begin{itemize}
  \item A single seed drives Python's \code{random}, NumPy, every estimator and
        the split.
  \item No paths, feature lists, ratios or hyperparameters are hardcoded; all
        live in \code{config.yaml}, validated on load.
  \item Every run writes a manifest recording hyperparameters, metrics on all
        four splits, feature set, seed, training time, model path, git commit
        hash and working-tree dirty flag.
  \item Dependencies are pinned in \code{requirements.txt}.
  \item The entire \code{experiments/} tree is regenerable with a single
        command.
\end{itemize}

\subsection{Review outcomes}

A strict pre-release review produced the following changes.

\begin{longtable}{@{}p{7.4cm} p{2.0cm} p{5.2cm}@{}}
\toprule
\textbf{Issue} & \textbf{Severity} & \textbf{Resolution} \\
\midrule
\endfirsthead
\toprule
\textbf{Issue} & \textbf{Severity} & \textbf{Resolution} \\
\midrule
\endhead
Imputation fit on full population       & Critical & fit/transform contract; regression tests \\
Train metrics never computed            & High     & all four splits evaluated; overfit warning \\
Results table re-read from disk         & Medium   & metrics carried in memory \\
Feature matrix re-converted per split   & Medium   & converted once \\
\code{assert} used for correctness      & Medium   & replaced with a typed exception \\
Model order depended on dict order      & Medium   & explicit ordering, sorted by run id \\
\code{sys.path} manipulation in scripts & Medium   & proper package via \code{pyproject.toml} \\
No dependency pinning                   & Medium   & \code{requirements.txt} with exact versions \\
Plot function coupled to table schema   & Low      & takes label/value sequences \\
Duplicated checks and cache loading     & Low      & extracted to shared helpers \\
Legacy scripts duplicating pipeline     & Low      & removed (history in git) \\
Unused imports and variables            & Low      & removed \\
\bottomrule
\caption{Code review findings and resolutions.}
\label{tab:review}
\end{longtable}

\subsection{Scale}

Approximately 2{,}600 lines of Python across 34 modules, with 24 unit tests
covering the invariants that matter: leakage prevention, patient-level split
isolation, ensemble arithmetic and configuration validation. Twelve trained
models and 59 generated figures.

% =============================================================================
\section{Limitations}
\label{sec:limitations}
% =============================================================================

\begin{itemize}
  \item \textbf{Biosignal modality unavailable.} The candidate pool is 12
        models rather than the planned 16.
  \item \textbf{Overfitting.} Four models reach train AUC 1.000. Hyperparameter
        grids were kept small for a two-core environment and should be widened.
  \item \textbf{Single seed.} All results use seed 42. The $\approx 0.008$ AUC
        ensemble margin has not been shown to exceed split-to-split noise.
  \item \textbf{SHAP for XGBoost.} \code{shap.TreeExplainer} fails on this
        shap/XGBoost version pair when parsing XGBoost's serialised
        \code{base\_score}. The failure is caught and logged; feature importance
        is unaffected. This is a version-compatibility issue, not a defect in
        the pipeline.
  \item \textbf{No biospecimen modality.} The columns exist in the curated cut
        but carry 60--80\% baseline missingness, requiring a deliberate
        missing-data decision before inclusion.
\end{itemize}

% =============================================================================
\section{Recommendations for Month 2}
% =============================================================================

\subsection{Handover}

Month 1 provides per-model predicted probabilities for all four splits on disk,
an untouched calibration split of 211 patients, a single call that reconstructs
every trained run, and an aggregation hook that accepts an arbitrary weight
vector. Month 2 requires no retraining to begin.

\subsection{A theoretical caveat that should be settled early}

Magesh et al.\ detect whether \emph{one input} is out-of-distribution by
combining $K$ statistics from \emph{one model}. This project inverts that: it
detects whether \emph{one model} is an outlier by comparing $M$ models across
many inputs. The machinery --- conformal $p$-values, step-up multiple testing,
false-alarm control --- transfers structurally, but the guarantee does not
transfer automatically.

Their Theorem 2 bounds the false-alarm rate for a single test input against a
calibration set of \emph{exchangeable} inputs. In the per-model setting the
objects under test are models, and twelve models trained on overlapping data are
neither independent nor exchangeable in the required sense.

Two defensible outcomes exist, and which one holds should be established by
week~6 rather than deferred to the writing phase:

\begin{enumerate}
  \item \textbf{The guarantee transfers with modification}, for example by
        defining exchangeability over \emph{patients} rather than models: for
        each calibration patient, compute each model's deviation from the
        leave-one-model-out consensus and treat those per-patient deviations as
        the exchangeable units. This is the more promising route.
  \item \textbf{It does not transfer}, in which case the contribution is an
        empirically effective, principled-by-construction selection procedure
        benchmarked against the Month 1 baselines. This remains publishable;
        claiming a guarantee that does not hold would not be.
\end{enumerate}

\subsection{Suggested schedule}

\begin{description}[leftmargin=2.2cm,style=nextline]
  \item[Weeks 5--6] Define and compare at least three model-level outlier
    scores: deviation from the leave-one-model-out consensus, mean pairwise
    disagreement, and calibration error. Sanity-check against a known weak
    model --- clinical + NeuralNet, test AUC 0.707 --- which any working score
    should flag.
  \item[Weeks 6--7] Compute conformal $p$-values on the calibration split,
    bootstrapping over patients. Verify empirically that $p$-values are
    approximately uniform under the null; this is where a broken
    exchangeability assumption will reveal itself.
  \item[Weeks 7--8] Apply the step-up procedure. Because model scores are
    strongly dependent, use Benjamini--Yekutieli rather than plain
    Benjamini--Hochberg. Aggregate survivors by trimmed mean or
    confidence-weighted average, and ablate removal versus down-weighting,
    the level $\alpha$, and Bonferroni as a comparison.
\end{description}

\subsection{Principal risks}

\begin{itemize}
  \item \textbf{The margin is thin.} A 0.008 AUC difference on 142 test patients
        is within plausible split noise. Run at least five seeds and report mean
        and standard deviation; lead with variance across conditions rather than
        mean performance, which is the claim the method is designed to support.
  \item \textbf{Pruning may not help when all models are competent.} Construct
        conditions where genuine outlier models exist --- the five
        distribution-shift datasets already assembled --- and verify that a
        deliberately degraded model (trained on shuffled labels) is removed. A
        method that cannot remove a label-shuffled model is not working,
        whatever its AUC.
  \item \textbf{Overfit models will distort scores.} Address regularisation
        before building selection logic on their validation behaviour.
  \item \textbf{Test-set discipline.} Every design decision must be made on the
        validation and calibration splits. The current code enforces this by
        convention rather than by construction.
\end{itemize}

% =============================================================================
\section{Conclusion}
% =============================================================================

Month 1 delivered a reproducible, configuration-driven baseline pipeline with a
12-model candidate pool and five ensemble baselines, together with the
documentation, diagrams and tests needed for another researcher to reproduce and
extend it. The most consequential outcomes were not the headline metrics but the
two leakage problems identified and corrected: without those fixes, the reported
test scores would have been optimistic in a way that would likely have survived
into a submitted paper.

The baseline that Month 2 must beat is 0.9894 test ROC-AUC from simple
averaging. That bar is high and the available margin is narrow, which is itself
a useful finding to have established before the research contribution is built
rather than after.

\end{document}
